\section{Ứng dụng mô hình thực tế}

\subsection{Ứng dụng 1: Tác động Trực tiếp Trung bình (Average Direct Effect --- ADE)}

\subsubsection{Mô hình và đại lượng mục tiêu}

\textbf{Bối cảnh:} Dữ liệu quan sát (observational data). Mục tiêu là cô lập tác động khi lật đối xử của \textit{một} đơn vị, tránh sai lệch do tương tác chéo (SUTVA vi phạm), trong khi giữ nguyên phân phối đối xử của các đơn vị khác.

Mô hình kết quả trong bài toán ADE:
\[
    Y_t = f^*(D_t, X_t, H_t) + \tilde{Y}_t, \qquad \mathbb{E}[\tilde{Y}_t \mid D_t, X_t, H_t] = 0
\]

\textbf{Đại lượng mục tiêu (ADE):}
\[
    \psi^* = \frac{1}{T}\sum_{t=1}^T \mathbb{E}\bigl[Y_t(1, H_t) - Y_t(0, H_t)\bigr]
\]

Đây là tác động trực tiếp trung bình khi thay đổi đối xử của một đơn vị từ 0 sang 1, giữ nguyên trạng thái chung \(H_t\) theo tiến trình tiến hóa tự nhiên của hệ thống. Kỳ vọng \(\mathbb{E}\) lấy trung bình theo phân phối \(W_{1:T}\), tức đã được gạt bỏ (marginalized) theo các giá trị thực của \(H_t\).

\subsubsection{Bộ ước lượng ADE (AIPW có điều chỉnh \(H_t\))}

Bộ ước lượng được xây dựng dưới dạng trung bình mẫu của hàm score \(\phi\):
\[
    \hat{\psi}_{\text{ADE}} = \frac{1}{T} \sum_{t=1}^T \phi(W_t;\, \eta)
\]

Hàm score \(\phi\) gồm hai thành phần:
\begin{align}
    \phi(W_t;\eta) &= \underbrace{\left[ f(1,X_t,H_t) - f(0,X_t,H_t) \right]}_{\text{Plug-in: dự đoán chênh lệch kết quả}} \nonumber \\
    &\quad + \underbrace{\left[ \left(\frac{D_t}{m(X_t)} - \frac{1-D_t}{1-m(X_t)}\right) \bigl(Y_t - f(D_t,X_t,H_t)\bigr) \right]}_{\text{Debiasing: sửa sai bằng propensity score}}
\end{align}

\begin{itemize}
    \item \textbf{Plug-in:} Dùng học máy dự đoán trực tiếp chênh lệch kết quả \(f(1,\cdot) - f(0,\cdot)\).
    \item \textbf{Debiasing:} Dùng propensity score \(m(X_t)\) để đánh trọng số lại phần dư và sửa lỗi cho mô hình kết quả. Điểm khác biệt so với AIPW chuẩn: mô hình kết quả \(f\) \textbf{bắt buộc phải nhận \(H_t\) làm đầu vào} để kiểm soát confounding qua trạng thái chung.
\end{itemize}

Dưới Assumption 4.1 (regularity: propensity bị chặn, outcome bounded) và Assumption 4.2 (tốc độ hội tụ: \(\hat{m}, \hat{f}\) hội tụ với tích tốc độ \(o_P(T^{-1/2})\)), \textbf{Theorem 4.3} kết luận:
\[
    \sqrt{T}\, \hat{\sigma}^{-1} \bigl(\hat{\psi}_{\text{ADE}} - \psi^*\bigr) \xrightarrow{d} \mathcal{N}(0,\, 1)
\]

\subsection{Ứng dụng 2: Tác động Toàn cục Trung bình (GATE) trong Switchback Experiments}

\subsubsection{Mô hình và đại lượng mục tiêu}

\textbf{Bối cảnh:} Thử nghiệm Switchback --- người thực hiện thí nghiệm kiểm soát toàn bộ \(D_{1:T}\) theo các khối thời gian kế tiếp (bật/tắt). Mục tiêu là so sánh khi \textit{tất cả} đơn vị được đối xử so với khi \textit{tất cả} bị kiểm soát, bao gồm cả tác động lan tỏa.

Mô hình kết quả trong Switchback với \(m\)-dependence:
\begin{align}
    Y_t(D_t, H_t) &= f^*(D_t, X_t, H_t) + \tilde{Y}_t \\
    H_t(D_{(t-m):(t-1)}) &= \operatorname{vec}(D_{(t-m):(t-1)},\, X_{(t-m):(t-1)})
\end{align}

\textbf{Đại lượng mục tiêu (GATE):}
\[
    \psi^* = \frac{1}{T}\sum_{t=1}^T \mathbb{E}\bigl[Y_t(1, H_t(\mathbf{1})) - Y_t(0, H_t(\mathbf{0}))\bigr]
\]

Đây là tác động toàn hệ thống (so sánh hai trạng thái ổn định: toàn treatment vs.\ toàn control), bao gồm cả hiệu ứng lan tỏa qua trạng thái chung.

\subsubsection{Bộ ước lượng GATE (Regression Adjustment + Debiasing)}

Hàm score cho GATE:
\begin{align}
    \phi(W_t;\eta) &= \underbrace{f(1,X_t,H_t(\mathbf{1})) - f(0,X_t,H_t(\mathbf{0}))}_{\text{Plug-in (Regression Adjustment)}} \nonumber \\
    &\quad + \underbrace{\Bigl(\frac{\mathbf{1}\{D_{t-m:t}=\mathbf{1}\}}{\pi^*(\cdot,1)} - \frac{\mathbf{1}\{D_{t-m:t}=\mathbf{0}\}}{\pi^*(\cdot,0)}\Bigr) \cdot (Y_t - f(D_t,X_t,H_t))}_{\text{Debiasing (chỉ kích hoạt ở cửa sổ thuần nhất)}}
\end{align}

\begin{itemize}
    \item \textbf{Plug-in (Regression Adjustment):} Dùng học máy dự đoán kết quả tiềm năng, hoạt động ở \textit{mọi} bước thời gian --- đây là cơ chế giúp \textbf{giảm đáng kể phương sai}.
    \item \textbf{Debiasing:} Chỉ kích hoạt khi chuỗi đối xử đạt trạng thái \textbf{thuần nhất} trong ít nhất \(m+1\) bước liên tiếp (cùng tất cả là đối xử hoặc kiểm soát). Điều này cho phép tận dụng toàn bộ dữ liệu thay vì bỏ đi dữ liệu ở giai đoạn chuyển giao.
\end{itemize}

Dưới Assumption 5.1 (xác suất chuỗi thuần nhất bị chặn: \(\pi^*((t{-}m):t, d) \geq \zeta > 0\)) và điều kiện tốc độ hội tụ tương tự, bộ ước lượng GATE cũng đạt chuẩn tiệm cận bình thường.
